\chapter{Introduction}

\section{What Is Connector?}

Connector is an \textbf{infrastructure-grade agentic runtime}: a platform that executes
autonomous AI agents the way an operating system executes processes — with governed
memory, typed communication channels, resource budgets, audit trails, and
cryptographic proof of every state transition.

Most AI agent frameworks are \emph{libraries}: they give developers helpful
abstractions for calling LLMs and chaining tool calls, but they leave
governance, auditability, recovery, and cost control entirely to the
application layer. In production environments — healthcare, finance, legal,
enterprise — this is insufficient. When an agent makes a decision, operators
need to know \emph{what} it decided, \emph{why}, \emph{what data it accessed},
\emph{how much it cost}, and \emph{whether a human can verify it} after the fact.

Connector solves this at the infrastructure level, not the application level.

\begin{conceptbox}{Core Value Proposition}
Every agent decision on Connector is:
\begin{itemize}[noitemsep]
  \item \textbf{Recorded} — written to a signed, content-addressed memory packet
  \item \textbf{Auditable} — every access, tool call, and state transition is logged
  \item \textbf{Traceable} — a full decision chain can be reconstructed at any time
  \item \textbf{Provable} — cryptographic receipts allow independent verification
  \item \textbf{Governed} — budgets, predicates, and compliance rules are enforced by the kernel
\end{itemize}
\end{conceptbox}

\section{The Problem Space}

Running AI agents in production creates a set of hard infrastructure problems
that cannot be solved by wrapping LLM APIs:

\begin{enumerate}
  \item \textbf{State fragility} — Agent state lives in memory and is lost on crash.
        Connector persists all state as immutable memory packets with lineage tracking.

  \item \textbf{No audit trail} — Log lines are not evidence. Connector writes
        cryptographically signed receipts for every operation.

  \item \textbf{Unbounded cost} — LLM token spend is invisible until the bill arrives.
        Connector enforces per-session and per-agent token budgets in real time.

  \item \textbf{No governance} — Compliance requirements (HIPAA, SOC2, GDPR) cannot
        be bolted on after the fact. Connector treats predicates and compliance
        rules as first-class kernel primitives.

  \item \textbf{Protocol chaos} — Agents need to communicate with each other, with
        external tools, and with humans through different transports. Connector
        provides a unified protocol stack (CNP) with bridges to A2A and MCP.
\end{enumerate}

\section{Architecture Overview}

Connector is structured as seven discrete layers, each building on the one below.
The architecture is intentionally kernel-inspired: lower layers provide
stable primitives; upper layers compose those primitives into product surfaces.

\begin{figure}[H]
\centering
\begin{tikzpicture}[node distance=0.55cm]

  \node[layerbox=CInfra] (infra)
    {\textbf{Layer 7 — Infrastructure}\quad
     Boot Sequence \textbullet\ Health Probes \textbullet\ Monitoring \textbullet\ Deployment};

  \node[layerbox=CBiz, below=of infra] (biz)
    {\textbf{Layer 6 — Business Operations}\quad
     Signup \textbullet\ Login \textbullet\ API Keys \textbullet\ Licensing \textbullet\ Billing};

  \node[layerbox=CDX, below=of biz] (dx)
    {\textbf{Layer 5 — Developer Experience}\quad
     Glue Facade \textbullet\ connectorctl CLI \textbullet\ REST API \textbullet\ Rust SDK};

  \node[layerbox=CLogic, below=of dx] (logic)
    {\textbf{Layer 4 — Logic Plane}\quad
     Contract Registry \textbullet\ Composition Patterns \textbullet\ Validation};

  \node[layerbox=CProto, below=of logic] (proto)
    {\textbf{Layer 3 — Protocol Layer}\quad
     CNP (7-layer native) \textbullet\ A2A Bridge \textbullet\ MCP Bridge};

  \node[layerbox=CExec, below=of proto] (exec)
    {\textbf{Layer 2 — Execution Engines}\quad
     Cognitive Substrate (11-layer) \textbullet\ CLS Compiler \textbullet\ VM Runtime};

  \node[layerbox=CKernel, below=of exec] (kernel)
    {\textbf{Layer 1 — Memory Kernel}\quad
     MemPacket Store \textbullet\ Syscall Interface \textbullet\ Agent Lifecycle \textbullet\ Predicates};

  \foreach \a/\b in {kernel/exec, exec/proto, proto/logic, logic/dx, dx/biz, biz/infra}{
    \draw[arr=gray!60] (\a.north) -- (\b.south);
  }

  %% Side annotations
  \node[right=0.4cm of kernel, font=\tiny\itshape, text=CKernel, align=left]
    {Foundation\\(immutable state)};
  \node[right=0.4cm of infra, font=\tiny\itshape, text=CInfra, align=left]
    {Operations\\(production-ready)};

\end{tikzpicture}
\caption{Connector 7-layer architecture — each layer provides stable primitives to the layer above}
\label{fig:arch-stack}
\end{figure}

The critical insight is \textbf{directionality}: requests flow \emph{down} through
the stack, and evidence flows \emph{up}. A developer calls a Glue verb; the
execution engine interprets it; the kernel enforces governance and writes a
signed receipt; that receipt bubbles back up through the stack as proof.

\section{Design Principles}

Seven principles govern every architectural decision in Connector:

\begin{table}[H]
\centering
\renewcommand{\arraystretch}{1.4}
\begin{tabularx}{\textwidth}{@{}lX@{}}
\toprule
\textbf{Principle} & \textbf{What it means in practice} \\
\midrule
\textbf{Memory-First} &
  All state is a MemPacket. Computation is memory transformation.
  Nothing is ephemeral unless explicitly marked as such. \\
\textbf{Content Addressing} &
  Every artifact is identified by its content hash (CID).
  Identical content is deduplicated; any tampering changes the CID. \\
\textbf{Governance by Default} &
  Predicates, budgets, and audit requirements are enforced by the kernel
  before any agent operation executes — not after. \\
\textbf{Recovery by Default} &
  Every layer checkpoints state. Node crash, budget violation, or network
  partition never causes unrecoverable data loss. \\
\textbf{Evidence Everywhere} &
  Every state transition produces a cryptographic receipt.
  Auditors can verify the full history of any agent at any time. \\
\textbf{Protocol Agnostic} &
  CNP is native, but A2A and MCP bridges mean Connector agents
  interoperate with the broader AI ecosystem out of the box. \\
\textbf{Operator Ergonomics} &
  The Glue facade and \texttt{connectorctl} CLI expose complex primitives
  through a simple, human-readable grammar. \\
\bottomrule
\end{tabularx}
\caption{Connector's seven design principles}
\end{table}

\section{Deployment Models}

Connector supports four deployment topologies:

\begin{figure}[H]
\centering
\begin{tikzpicture}[
  node distance=1.4cm and 1.8cm,
  model/.style={rectangle, rounded corners=5pt, draw=MidnightBlue!60,
                fill=MidnightBlue!8, text width=3.2cm, align=center,
                minimum height=1.6cm, font=\small},
  lbl/.style={font=\footnotesize\itshape, text=gray}
]

\node[model] (local)   {Local Node\\[4pt]\footnotesize Single machine\\Dev \& edge};
\node[model, right=of local]  (cluster) {Cluster\\[4pt]\footnotesize Multi-node\\Distributed memory};
\node[model, right=of cluster] (saas)   {Hosted SaaS\\[4pt]\footnotesize Fully managed\\Cloud control plane};
\node[model, right=of saas]   (hybrid) {Hybrid\\[4pt]\footnotesize Local execution\\Cloud licensing};

\draw[arr=MidnightBlue!40] (local) -- node[above,lbl]{scale up} (cluster);
\draw[arr=MidnightBlue!40] (cluster) -- node[above,lbl]{delegate} (saas);
\draw[darr=CProto!60, dashed] (local) -- node[below,lbl]{license sync} (hybrid);

\end{tikzpicture}
\caption{Connector deployment topology options}
\end{figure}

\begin{itemize}[leftmargin=1.5em]
  \item \textbf{Local Node} — A single binary (\texttt{connector-platform}) runs the entire
        stack on one machine. Used for development, testing, and edge deployments.
  \item \textbf{Cluster} — Multiple nodes share a distributed memory store with
        consensus for high availability and horizontal scale.
  \item \textbf{Hosted SaaS} — The Connector team operates the control plane;
        customers access it via the customer portal and API keys.
  \item \textbf{Hybrid} — Execution stays on-premise (data sovereignty), while
        licensing, billing, and the admin dashboard run on the hosted control plane.
\end{itemize}

\section{Key Terminology}

\begin{defbox}{Glossary}
\begin{description}[leftmargin=2.4cm, style=nextline, itemsep=4pt]
  \item[\texttt{MemPacket}]
    Immutable, content-addressed memory object. The atomic unit of all state in Connector.
    Every MemPacket has a CID, typed content, metadata (timestamps, lineage), and an
    optional Ed25519 signature.
  \item[\texttt{CID}]
    Content Identifier — \texttt{cls1-sha256-<64-char-hex>}. Computed as the SHA-256
    hash of the canonical JSON representation of the packet's content. Any change
    to content produces a different CID.
  \item[\texttt{Agent}]
    Autonomous execution unit with a DID (decentralised identifier), memory context,
    capability grants, and a lifecycle state machine (Registered → Active → Paused →
    Terminated).
  \item[\texttt{Session}]
    Scoped execution context. Carries a resource budget (token limit, memory quota,
    wall-clock limit), a set of governance predicates, and a checkpoint chain.
  \item[\texttt{Contract}]
    Compiled, signed executable logic produced by the CLS compiler. Contains an
    interface specification, a state machine, governance rules, and a resource envelope.
  \item[\texttt{CNP}]
    Connector Native Protocol — a 7-layer typed communication stack for
    agent-to-agent messaging with capability gating, encryption, and delivery guarantees.
  \item[\texttt{CLS}]
    Connector Logic Standard — the contract language, 6-pass compiler, and runtime
    VM that transform agent intent declarations into governed, verifiable execution.
  \item[\texttt{Glue}]
    The developer-facing surface of Connector. Provides action-first verbs
    (\texttt{start}, \texttt{stop}, \texttt{list}, \texttt{trace}, \texttt{prove}) over
    CLI, REST API, and Rust library interfaces.
  \item[\texttt{SOE}]
    Surface Output Engine — the layer that produces signed, CID-anchored output
    documents (reports, receipts, defense packages) for auditors and operators.
\end{description}
\end{defbox}

\section{How to Read This Document}

This document progresses \emph{bottom-up} through the architecture stack,
mirroring how the platform is built. Readers who want to understand how an
operator deploys and operates Connector should read Chapters 6--7 first.
Readers building agents or integrations should focus on Chapters 2--5.

\begin{table}[H]
\centering
\renewcommand{\arraystretch}{1.3}
\begin{tabular}{@{}clp{7.5cm}@{}}
\toprule
\textbf{Ch.} & \textbf{Layer} & \textbf{Covers} \\
\midrule
2 & Memory Kernel    & MemPacket anatomy, CID addressing, syscall categories,
                       predicate enforcement, agent lifecycle, cryptographic receipts \\
3 & Execution Engines & 11-layer Cognitive Substrate, CLS 6-pass compiler,
                       contract runtime, resource budgets \\
4 & Protocol Layers  & CNP 7-layer stack, A2A and MCP bridges, message types,
                       capability attenuation, delivery guarantees \\
5 & Developer UX     & Glue verb grammar, connectorctl CLI, REST API surface,
                       configuration hierarchy, error model \\
6 & Business Ops     & Two-server architecture, signup/login flows, API key
                       issuance, license tiers, Stripe billing integration \\
7 & Node Operations  & 12-stage boot sequence, health probes, agent deployment
                       manifest, monitoring, disaster recovery \\
\bottomrule
\end{tabular}
\caption{Chapter guide by audience and layer}
\end{table}
